% from ega4-14-fr.tex
\oldpage[IV-3]{212}

Lorsque de plus $Y$ est localement noethérien, ces conditions sont encore
équivalentes à la suivante~:
\begin{enumerate}[label=\emph{\alph*)}]
  \item[c)] $f$ est ouvert en tous les points de $X_y$ (ou, ce qui revient au même
  (14.3.3.1, (ii)), aux points maximaux de $X_y$).
\end{enumerate}
\end{corollary}

L'équivalence de $a)$ et $c)$ lorsque $Y$ est localement noethérien résulte de
(14.4.1)~; les conditions $b)$ ou $b')$, appliquées aux \emph{points maximaux} de
$X_y$, entraînent $a)$ en vertu également de (14.3.3.1, (ii))~; enfin, $b')$ et
$b'')$ sont équivalentes, car
\[
  \dim(U\cap X_y)=\sup_{x\in X_y}(\dim_x(U\cap X_y)).
\]
Reste à voir que la condition $a)$ entraîne $b)$ et $b')$ en \emph{tout} point
$x\in X_y$. Posons $d=\dim_x(X_y)$, et soit $x'$ le point générique d'une
composante irréductible de $X_y$ contenant $x$ et de dimension $d$. En vertu de
$a)$ et de (14.4.1), il y a une composante irréductible $Z$ de $X$ contenant $x'$
et équidimensionnelle sur $Y$ au point $x'$, donc telle que
$\dim_{x'}(Z\cap X_y)=\dim(Z\cap X_\eta)$. Mais par construction
$\dim_{x'}(Z\cap X_y)=\dim_{x'}(X_y)=d$, et
$\dim_x(Z\cap X_y)\leq\dim(Z\cap X_y)=d$~; compte tenu de (13.1.6), cela prouve
que $Z$ est équidimensionnel sur $Y$ au point $x$~; donc $a)$ entraîne $b)$. De
plus, on a $\dim(X_\eta)\geq\dim(Z\cap X_\eta)=d=\dim_x(X_y)$. Remplaçant $X$ par
un voisinage ouvert $U$ de $x$, on voit ainsi que $a)$ entraîne $b')$. C.Q.F.D.

\begin{corollary}[14.4.3]
\label{IV.14.4.3-fr}
Soient $Y$ un préschéma géométriquement unibranche, $f:X\to Y$ un morphisme
localement de type fini. Les conditions suivantes sont équivalentes~:
\begin{enumerate}[label=\emph{\alph*)}]
  \item $f$ est universellement ouvert.
  \item Pour tout ouvert $U$ de $X$, tout $y\in Y$ et toute générisation $y'$ de
  $y$, on a $\dim(U\cap X_{y'})\geq\dim(U\cap X_y)$.
\end{enumerate}
Si de plus $Y$ est localement noethérien, ces conditions sont aussi équivalentes à
la suivante~:
\begin{enumerate}[label=\emph{\alph*)}]
  \item[c)] $f$ est ouvert.
\end{enumerate}
\end{corollary}

\begin{corollary}[14.4.4]
\EGAReaderTarget{ega4.c04.cor.14-4-4}{IV.14.4.4-fr}
(critère de Chevalley). Soit $f:X\to Y$ un morphisme localement de type fini.
\begin{enumerate}[label=(\roman*)]
  \item Si $f$ est équidimensionnel en un point $x\in X$ (13.3.2) et si
  $y=f(x)$ est un point géométriquement unibranche de $Y$, $f$ est
  universellement ouvert au point $x$.
  \item Si $Y$ est géométriquement unibranche, $f$ est universellement ouvert en
  tous les points de $X$ où $f$ est équidimensionnel, et l'ensemble de ces points
  est ouvert dans $X$. En particulier, si $f$ est équidimensionnel, il est
  universellement ouvert.
\end{enumerate}
\end{corollary}

L'assertion (ii) est conséquence triviale de (i), puisqu'on sait déjà que l'ensemble
des points où $f$ est équidimensionnel est ouvert (13.3.2). Quant à l'assertion (i),
elle résulte de ce que l'hypothèse implique que la condition $b)$ de (14.4.1) est
vérifiée au point générique d'une composante irréductible de $X_y$ contenant $x$,
compte tenu de (13.3.1)~; il suffit donc d'appliquer (14.4.1).

\begin{remark}[14.4.5]
\label{IV.14.4.5-fr}
On peut prouver que si $Y$ est localement noethérien, et si toutes les générisations
de $y=f(x)$ sont des points géométriquement unibranches de $Y$ (cf. (6.15.2)),
alors, si $f$ est équidimensionnel au point $x$, il est universellement ouvert
\emph{dans un voisinage de $x$}.
\end{remark}

